% BHU PYQ -- Manually transcribed & error-corrected
% Paper: BCH-322 : Goods and Service Tax Act and Customs Acts
% Exam: B.Com. (Hons.) Semester VI Examination, 2022-23

\documentclass[11pt,a4paper]{article}
\usepackage[a4paper,top=2.5cm,bottom=2.5cm,left=2.8cm,right=2.8cm,headheight=14pt]{geometry}
\usepackage{amsmath,amssymb}
\usepackage[T1]{fontenc}
\usepackage[utf8]{inputenc}
\usepackage{lmodern,microtype,fancyhdr,enumitem,hyperref,lastpage,parskip}

\hypersetup{
    colorlinks=true,
    urlcolor=black,
    linkcolor=black,
    pdftitle={BCH-322: Goods and Service Tax Act and Customs Acts -- BHU B.Com. (Hons.) Sem VI 2022-23}
}

\pagestyle{fancy}
\fancyhf{}
\renewcommand{\headrulewidth}{0.4pt}
\renewcommand{\footrulewidth}{0.4pt}
\fancyhead[L]{\small   \textit{Banaras Hindu University}}
\fancyhead[R]{\small   \textit{BCH-322: Goods and Service Tax Act and Customs Acts}}
\fancyfoot[C]{\small Page  \thepage\ of \pageref{LastPage}}


\newlist{parts}{enumerate}{2}
\setlist[parts,1]{label=(\alph*),leftmargin=*,itemsep=5pt,topsep=3pt}
\setlist[parts,2]{label=(\roman*),leftmargin=*,itemsep=3pt,topsep=2pt}

\newcommand{\pts}[1]{\hfill   \textit{[#1]}}


\begin{document}


\begin{center}
  {\large\bfseries BANARAS HINDU UNIVERSITY}\\[2pt]
  {
\normalsize B.Com. (Hons.) --- Semester VI Examination, 2022-23}\\[6pt]
  \rule{0.8\linewidth}{0.5pt}\\[6pt]
  {\large\bfseries Paper No. BCH-322: Goods and Service Tax Act and Customs Acts}\\[4pt]
  {
\normalsize Time: 3 Hours\hfill Maximum Marks: 70}
\end{center}

\rule{\linewidth}{0.4pt}

\smallskip



\noindent   \textit{Instructions: Attempt all questions. Each question carries equal marks (14 marks each).}

\bigskip




\noindent   \textbf{Question 1.}
Explain the reasons for adopting GST in India and its merits and limitations. Explain the tax laws that got replaced by GST.\pts{14}

\centerline{   \textbf{OR}}

Write a detailed note on exempted supply, non-taxable supply, zero rated supply and export of goods and services.\pts{14}

\medskip\rule{\linewidth}{0.2pt}




\noindent   \textbf{Question 2.}
What is mandatory registration? How does it differ from voluntary registration? Explain the different types of registration in GST.\pts{14}

\centerline{   \textbf{OR}}

What details are furnished in GSTR-1 and GSTR-2? What are the procedure and time limit applicable for the revision/rectification of details furnished in GSTR-1 and GSTR-2?\pts{14}

\medskip\rule{\linewidth}{0.2pt}




\noindent   \textbf{Question 3.}
What items are included and excluded in transaction value for the computation of taxable value under GST?\pts{14}

\centerline{   \textbf{OR}}

XYZ Ltd., a registered manufacturer in the state of Uttar Pradesh, provides the following particulars for the tax period of April, 2022:

\begin{parts}
  \item Inputs purchased within the state: Rs. 5,04,000 (including 12\% GST).
  \item Machinery purchased on 12-04-2022 for Rs. 3,00,000 (excluding 12\% GST) from a dealer in Delhi is eligible for the input tax credit. Depreciation is charged at 10\% p.a.
  \item Manufacturing expenses: Rs. 1,60,000.
  \item The profit margin is 20\% of the sales price.
  \item The goods produced were sold to a buyer in Patna with 18\% IGST on sales.
  \item A 5\% discount was allowed at the time of sale, and subsequently a 2\% discount was allowed at the time of payment.
\end{parts}
Calculate the amount of GST payable after utilizing the input tax credit for the month of April, 2022, assuming that no opening balance of input tax credit was available.\pts{14}

\medskip\rule{\linewidth}{0.2pt}




\noindent   \textbf{Question 4.}
Describe the features of Customs Duty in India. Explain the role of Customs Duty in International Trade.\pts{14}

\centerline{   \textbf{OR}}

What do you mean by the import of Goods? Describe the various categories under which imports can be made.\pts{14}

\medskip\rule{\linewidth}{0.2pt}




\noindent   \textbf{Question 5.}
Elucidate the valuation rules under Customs Duty. Describe items includable in determining the Assessable Value under Customs Act.\pts{14}

\centerline{   \textbf{OR}}

An Indian importer has imported raw material for 5,000 dollars in August, 2022. Following information is available:

\begin{parts}
  \item Packaging charges of goods: 120 dollars.
  \item Goods were stuffed in a container (returnable); price of the container is 400 dollars.
  \item Insurance premium: 50 dollars.
  \item Sea freight: 160 dollars.
  \item The importer had paid a commission of 100 dollars to a broker who arranged the transaction.
  \item Dollar Rate is Rs. 80 = 1 dollar.
  \item Basic Customs Duty is 10\%.
  \item Integrated tax u/s 3(7) of Customs Tariff Act 1975 is 12\%.
\end{parts}
Find out the assessable value of imported goods and Customs Duty Payable.\pts{14}

\vfill
\rule{\linewidth}{0.4pt}

\begin{center}\small
\href{https://bhu-exam.vercel.app/}{Click here to visit our website}\\[4pt]
\href{https://me-aryan.vercel.app/}{Click here to visit our contributor}\\[4pt]
\href{https://quashan-me.vercel.app/}{Click here to visit our contributor}
\end{center}

\end{document}
